\section{A direct map from young forcing to the old stress}
\label{sec:young-old-direct}

This section turns the actual train's forward estimate and the
all-age outputs of Section~\ref{sec:canonical-transfer} into a
uniform map on young forcing. The physical pressure and curved-space
comparison are treated separately.

Put $a(s)=1+\nu s$, $z(s)=s/a(s)$,
$\theta(s)=1/[\nu a(s)]$, $e=(\nu c^2)^{-1}$, and $c=\nu^{-2}$.
Fix $K\Subset\{0<\operatorname{Re}\mu<1/4\}$, $1<\alpha<1.2$,
and a fixed $\Sigma>\alpha$.
Write $H=\Sigma/\nu$, $T_f=\alpha/\nu$, and
$u_o=\nu^{-3/2}$. The auxiliary time $s$ is distinct from physical
logarithmic time. All cell norms use the unnormalized spatial measure.
The weighted vorticity space is the centrally even, zero-circulation
subspace $X_{m,c}$, with a fixed $m>8$; its first moments vanish.

Let $\mathcal U(s,t)$ be the actual train evolution, and let
$q_{\pm,s}$ and $\ell_{\pm,s}$ be the heat-scaled, periodized frozen
Gaussian pair and its finite Gram-normalized dual. This Gram normalization
only makes $\ell_{i,s}q_{j,s}=\delta_{ij}$; it is not a terminal
measurement. Set
\[
 F_\pm(s,t)=
 \exp\!\left[\lambda_\pm(z(s)-z(t))\right]
 \left(\frac{a(s)}{a(t)}\right)^{\kappa_\pm},
 \quad(\lambda_+,\kappa_+)=(\lambda,\kappa),
 \quad(\lambda_-,\kappa_-)=(\bar\lambda,\bar\kappa).
\]
The two diagonal factors have equal modulus. On $0\le s\le H$,
the direct/dual residual is $\delta_c$ and the following are the
retained block estimates: relative complementary propagation is
$Ce^{-g(s-t)}$, the direct amplitude error is $C\delta_c$, and both
amplitude/complement couplings are $C(\nu+\delta_c)$. One has
\[
 \delta_c=C_\Sigma(c^{-2}+\delta_{\rm base}),\qquad
 \delta_{\rm base}=P_\nu e^{C_\Sigma/\nu-b_\Sigma c^2},\qquad
 \eta_\nu=\nu+\delta_c/\nu=O(\nu).
\]
These bounds follow from the block construction in
Proposition~\ref{prop:short-full-clock}: on $a\le1+\Sigma$,
$\theta^{-1}=\nu a$, $d^{-2}=e/\theta=O_\Sigma(c^{-2})$,
and conversion to physical auxiliary time divides the generator by
$a^2$. Its two diagonal entries are
$\lambda_\pm/a^2+\kappa_\pm\nu/a$. Moreover
$z(s)-z(t)\ge(s-t)/(1+\Sigma)^2$, and the endpoint heat-scaling
maps have uniform norms on this finite band. Thus the moving
complementary decay is $Ce^{-g_\Sigma(s-t)}$ and the stated
direct and coupling errors follow.

Divide the two amplitudes by $F_\pm(s,t)$ and the complement by
their common modulus. If $A$ and $R$ are the resulting amplitude
and complementary norms, variation of constants gives
\[
 R(s)\le Ce^{-g(s-t)}R(t)
       +C(\nu+\delta_c)\int_t^s e^{-g(s-r)}A(r)\,dr.
\]
Insert this estimate in the amplitude equation. Its deviation from
the initial amplitude is bounded by
\[
 C\left[\delta_c(s-t)+(\nu+\delta_c)
                 +(\nu+\delta_c)^2(s-t)\right]
       (|A(t)|+R(t))
\]
times a Gronwall factor bounded uniformly on $s-t\le\Sigma/\nu$.
Indeed the coefficient of that factor is
$C[\delta_c+(\nu+\delta_c)^2]$, whose integral is bounded and
tends to zero. Consequently, since $H-t\ge(\Sigma-\alpha)/\nu$ for
$0\le t\le T_f$,
\begin{equation}\label{eq:young-frozen-pair}
 \mathcal U(H,t)=
 \sum_{j=\pm}F_j(H,t)q_{j,H}\ell_{j,t}
       +|F_+(H,t)|\,O_{X_{m,c}\to X_{m,c}}(\eta_\nu).
\end{equation}
The exponentially decaying initial complement is absorbed in
$\eta_\nu$; the full error is not claimed exponentially small.
Keeping that complement in the same inequalities gives the full
gain bound on every subinterval.

Define the actual bare columns by
\[
 \phi_{\nu,+}(s)=\mathcal U(s,0)P_cq,\qquad
 \phi_{\nu,-}(s)=\overline{\phi_{\nu,+}(s)}.
\]
Using \eqref{eq:young-frozen-pair} at $t=0$ shows
\[
 \phi_{\nu,j}(H)=F_j(H,0)q_{j,H}
                          +|F_+(H,0)|O(\eta_\nu).
\]
Thus the same estimate may be written with these actual bare columns,
with coefficient $F_j(t,0)^{-1}\ell_{j,t}$ and the same relative
$O(\eta_\nu)$ error. This is an approximate covector derived directly
from the frozen pair; no exact transported observer is used.

\subsection{A uniform tested estimate for all later ages}

Fix the large clock $b$ used in the long-clock comparison.
The kinetic, tested-moment and absolute-stress impulse estimates
\eqref{eq:ct-kinetic-impulse}, \eqref{eq:ct-tested-impulse} and
\eqref{eq:ct-absolute-impulse} apply uniformly through overlap.
Only their value at $r=b$ is needed here.

For a state at $H$, \eqref{eq:short-causal-long-forward}
propagates it to $b$ with bound
\[
 \|a_b^2w(s_b,\sqrt{a_b}\,\cdot)/N_\nu\|_{X_{m,d(b)}}
   \le C_b |F_+(H,0)|^{-1}\|w(H)\|_{X_{m,c}},
 \qquad N_\nu=e^{\lambda/\nu}\nu^{-\kappa}.
\]
Above $b$ the same estimate has an exponentially decaying factor
$e^{-\gamma\theta}\theta^{-\operatorname{Re}\kappa}$.
At $b$, the uniform Biot--Savart and moment embedding applies since
$d(b)\ge1$. Equations \eqref{eq:ct-tested-impulse} and
\eqref{eq:ct-absolute-impulse} therefore show that every old
weighted moment and absolute stress functional used in the leading
strain calculation has norm at most
\begin{equation}\label{eq:young-later-output}
 C_K |F_+(H,0)|^{-1}\|w(H)\|_{X_{m,c}}
\end{equation}
after division by its canonical scale
$|N_\nu|\nu^{3+2\operatorname{Re}\mu}/c$.
No ordinary inverse factor or logarithmic loss occurs here.

\subsection{Young source theorem and angular signs}

Let $f$ be supported in $0<s\le T_f$ and valued in $X_{m,c}$.
Consider the zero-entrance causal equation
\[
 (\partial_s+\mu/s)w=L_{\nu,c}(s)w+f(s),
 \qquad s^\mu w(s)\longrightarrow0\quad(s\downarrow0).
\]
The tested source norm below defines a causal old-output theorem;
it alone does not imply a physical energy entrance. For the Gaussian
affine sources actually used here, $f(s)=O(s^{3/2})$ in the fixed
smooth weighted norms near zero. The local bounded forward evolution
then gives $w(s)=O(s^{5/2})$ by Duhamel, uniformly for $\mu\in K$.
Thus its planar velocity has finite entrance energy
$\int_0^1s^{-3/2}\|K_cw(s)\|_2^2ds<\infty$; the same fixed
derivative bounds give the required local derivative integrability.
The curved physical reconstruction still uses the separate free graph
and comparison theorem. A general source in the tested completion
must not be assigned these stronger properties without checking them.
Define the finite tested norm and explicit frozen covectors
\begin{align*}
 \|f\|_{\mathcal I_{\nu,\mu}}
  &=\int_0^{T_f}t^{\operatorname{Re}\mu}|F_+(t,0)|^{-1}
                         \|f(t)\|_{X_{m,c}}\,dt,\\
 a_{\nu,\pm}(f)
  &=\int_0^{T_f}t^\mu F_\pm(t,0)^{-1}\ell_{\pm,t}f(t)\,dt.
\end{align*}
Their bounds are viscosity-independent:
$|a_{\nu,+}|+|a_{\nu,-}|\le C_K\|f\|_{\mathcal I_{\nu,\mu}}$.
Since $a(t)\le1+\alpha$ on the source band, the norm is bounded by
$C_K\int_0^{T_f}t^{\operatorname{Re}\mu}e^{-\gamma t/(1+\alpha)}
\|f(t)\|\,dt$.
For a spectral scan, keep the source space fixed: choose a slightly
larger compact strip with endpoints $\sigma_-<\operatorname{Re}\mu
<\sigma_+$ and use $t^{\sigma_-}$ on $t\le1$ and
$t^{\sigma_+}$ on $t\ge1$ in the defining norm. The strict margins
absorb powers of $|\log t|$, so the displayed covectors and tested
outputs are holomorphic in $\mu$ for holomorphic source families.
All geometry, coefficients and cutoffs remain fixed during this scan.

Duhamel and \eqref{eq:young-frozen-pair} give exactly, at $H$,
\[
 H^\mu w(H)=
 a_{\nu,+}(f)\phi_{\nu,+}(H)
 +a_{\nu,-}(f)\phi_{\nu,-}(H)+r_H,
 \quad
 \|r_H\|\le C_K\eta_\nu|F_+(H,0)|
                                  \|f\|_{\mathcal I_{\nu,\mu}}.
\]
For $s\ge H$ both sides evolve by the same unforced equation.
Apply \eqref{eq:young-later-output}: the normalized old moment and
absolute stress errors are at most
$C_K\eta_\nu\|f\|_{\mathcal I_{\nu,\mu}}$.
This is an operator statement on the specified source class, not
only a calculation for one chosen forcing.

Put $J(\mu)=\mathcal J(\mu)$, $J_+(\mu)=J(\mu)$ and
$J_-(\mu)=\overline{J(\bar\mu)}$, and set
$\mathcal A_\nu=|N_\nu|\nu^{3+2\operatorname{Re}\mu}/c$.
The canonical Mellin limit and stress defect, with the opposite
angular identities of the two harmonics, give uniformly on the source class
\begin{align}
 I_T^\mu(w)
  &=-\frac{3\pi i}{4c}(\mu+3/2)\nu^{3+2\mu}
      \bigl[a_{\nu,+}(f)N_\nu J_+(\mu)
          -a_{\nu,-}(f)\overline{N_\nu}J_-(\mu)\bigr]
       +o_K(\mathcal A_\nu)\|f\|_{\mathcal I_{\nu,\mu}}.
       \label{eq:young-old-anisotropy}
\end{align}
Here $I_T^\mu$ is the integral of
$(2E_q-E_\eta)s^{-3/2}$, where, for the already Floquet-weighted
solution $w$ in this theorem,
\[
 E_\eta=\frac2c\int v_1(K_cw)_1,\qquad
 E_q=\frac2c\int v_2(K_cw)_2,\qquad
 C=\frac1c\int[v_1(K_cw)_2+v_2(K_cw)_1].
\]
These are bilinear pairings, without complex conjugation or another
factor $s^{-\mu}$.
For real $\mu$ and real forcing, $a_-=\bar a_+$, so the leading
term of \eqref{eq:young-old-anisotropy} is
$[3\pi/(2c)](\mu+3/2)\nu^{3+2\mu}
\operatorname{Im}(a_{\nu,+}N_\nu J(\mu))$.
The Floquet parameter is not conjugated in the negative channel
when it is complexified.

The source is zero before the old analytical tail begins, since
$T_f<u_o$. Thus the old moment law is homogeneous. Its lower endpoint
is exponentially small after division by $N_\nu$ and its upper
weighted endpoint vanishes; the old cutoff transition below $u_o$
is retained in its separate earlier-old contribution.
